\PassOptionsToPackage{colorlinks=true,allcolors=linkblue}{hyperref}
\documentclass[11pt,a4paper]{moderncv}
\moderncvstyle{banking}
\moderncvcolor{black}
\nopagenumbers{}

\usepackage[utf8]{inputenc}
\usepackage{newtxtext}
\usepackage{ragged2e}
%\usepackage[scale=0.92]{geometry}
\usepackage[
  a4paper,
  top=9mm,
  bottom=1mm,
  left=10mm,
  right=10mm
]{geometry}
\usepackage{multicol}
\usepackage{enumitem}
\usepackage{amssymb}
\setlist[itemize]{topsep=0pt, partopsep=0pt, parsep=0pt}
\definecolor{linkblue}{HTML}{1F3558}
\AtBeginDocument{%
  \hypersetup{colorlinks=true,allcolors=linkblue}%
}

\name{Leo}{Shaposhnik}

\newcommand*{\customcventry}[7][.13em]{
	\begin{tabular}{@{}l}
		{\bfseries #4}\\
		{\itshape #3}
	\end{tabular}
	\hfill
	\begin{tabular}{l@{}}
		{\bfseries #5}\\
		{\itshape #2}
	\end{tabular}
	\ifx&#7&%
	\else{\\
		\begin{minipage}{\maincolumnwidth}%
			\small#7%
	\end{minipage}}\fi%
	\par\addvspace{#1}}

\begin{document}
\makecvtitle
\vspace*{-13mm}

\begin{center}
	\large
	\begin{tabular}{ @{}c@{\hspace{10pt}}c@{\hspace{10pt}}c@{\hspace{10pt}}c@{\hspace{10pt}}c@{\hspace{10pt}}c@{\hspace{10pt}}c@{} }
		\href{https://replicawormhole.github.io/homepage/}{\faGlobe} &
		\href{mailto:leoshaposhnik@gmail.com}{\faEnvelope} &
		\href{https://github.com/ReplicaWormhole}{\faGithub} &
		\href{https://www.linkedin.com/in/leo-shaposhnik-7533801b5/}{\faLinkedin} &
		\href{https://orcid.org/0000-0001-8258-6233}{\faOrcid} &
		\href{https://inspirehep.net/authors/2605231}{\faAtom} \\
	\end{tabular}\normalsize
\end{center}

\vspace*{-8mm}

% =====================

% =====================
\section{Overview}
\vspace*{-2mm}
I am a theoretical physicist who builds agentic systems for open-ended scientific discovery.
In this capacity, I design and build adversarial AI systems for the development
of mathematical theories of physical systems and develop
code for the large-scale simulation of many-body systems on HPC clusters.
I have 10+ years of experience writing code, including 6+ years of experience
contributing to academic codebases and 1+ year contributing to production
codebases in industry as an IT consultant. I lead several research projects
from initial conception through completion and publication in international
journals.
Over the past year I have
developed a deep interest in knowledge representation, formal
verification and mechanistic
reasoning as approached by transformers, dependent type
theory and category theory.
\vspace*{-2mm}

\section{Current Position}
\vspace*{-2mm}

\customcventry{04/2023 -- present}
{Freie Universität Berlin}
{PhD in Theoretical Physics}
{Berlin, Germany}
{}{
\begin{itemize}[leftmargin=0.6cm, label=\textbullet]
	\setlength\itemsep{0.2em}
    \item Conduct research at the intersection of quantum information theory, tensor
        networks and high-energy physics.
    \item Experiment with multi-agent systems and context engineering for reliable
        collaboration on open-ended discovery.
    \item Implement parallelized high-performance code for the
        simulation of many-body systems on
        HPC clusters.
    \item Lead multiple international collaborations from the initial idea
        towards publication.
    \item Optimize RNNs and multiscale CNNs for the numerical study of quantum
        states.
    \item Co-organize a weekly seminar on quantum many-body physics.
	\item Teach courses on quantum mechanics, gravity and quantum information.
\end{itemize}
}
\section{Previous Experience}
\vspace*{-2mm}

\customcventry{11/2021 -- 03/2023}
{d-fine GmbH}
{IT Consultant}
{Berlin, Germany}
{}{
\begin{itemize}[leftmargin=0.6cm, label=\textbullet]
	\setlength\itemsep{0.2em}
	\item Owned a core software component of a production-grade IFRS 17
        reporting backend for an international reinsurer.
    \item Coordinated feature planning and implementation with
        client stakeholders.
    \item Supported testing, deployment and bug-fixing across the whole
        project.
\end{itemize}
}

\customcventry{07/2020 -- 10/2021}
{Physikalisch-Technische Bundesanstalt (PTB)}
{Working Student — Mathematical Modeling and Data Analysis}
{Berlin, Germany}
{}{
\begin{itemize}[leftmargin=0.6cm, label=\textbullet]
	\setlength\itemsep{0.2em}
	\item Formulated theoretical models to study roundoff-induced effects in classical chaotic systems.
		\item Performed high-precision long-time simulations of chaotic systems on
	        an HPC cluster.
\end{itemize}
}

\customcventry{\hspace{4pt} 03/2017 -- 06/2020}
{DESY}
{Working Student}
{Zeuthen, Germany}
{}{
\begin{itemize}[leftmargin=0.6cm, label=\textbullet]
	\setlength\itemsep{0.2em}
	\item Designed and set up optical experiments to measure
        radiation damage in optical components for a linear accelerator.
		\item Analyzed the concentration of radioactive material in accelerator
        components with MATLAB.
\end{itemize}
}

\section{Education}
\vspace*{-2mm}

\customcventry{10/2019 -- 02/2022}
{Humboldt-Universität zu Berlin / Albert Einstein Institute}
{Master’s Degree — Physics (Avg. Grade: 1.0 / GPA: 4.0)}
{Berlin, Germany}
{}{
Thesis on entanglement entropy in curved spacetime, including analytical derivations and numerical validation using tensor network techniques.
}

\customcventry{10/2016 -- 09/2019}
{Humboldt-Universität zu Berlin}
{Bachelor’s Degree — Physics (Avg. Grade: 1.7 / GPA: 3.3)}
{Berlin, Germany}
{}{
Thesis on Hopf-algebraic structures underlying perturbative renormalization in quantum field theory.
}
% =====================
\section{Technologies}
\vspace*{-2mm}
\setlength{\columnsep}{8pt}
\begin{multicols}{3}
\begin{itemize}[label=\textbullet, leftmargin=*, nosep, topsep=0pt, partopsep=0pt, parsep=0pt, itemsep=0pt]
    \item \textbf{Python}: PyTorch, Triton, SciPy, \newline Hugging Face Transformers
\item \textbf{High-performance computing}: C++, CUDA, Linux, Slurm
\item \textbf{Agent systems}: OpenAI API, Anthropic API, LangGraph, MCP, RAG
  \item \textbf{Data}: SQL, SAS, Pandas
  \item \textbf{Formal methods}: Lean 4
  \item \textbf{Scientific computing}: Julia, MATLAB, Wolfram Language
  \item \textbf{Version control \& CI/CD}: Git, Github, Azure DevOps


\end{itemize}
\end{multicols}

\section{Awards \& Honors}
\vspace*{-2mm}
\begin{itemize}[leftmargin=0.6cm, label=\textbullet]
	\setlength\itemsep{0.2em}
	\item Honorable mention for my poster ``On Infinite Tensor Networks, Complementary Recovery and Type II Factors''. QIQG 2025, Perimeter Institute.
\end{itemize}
\newpage
% =====================
\section{Publications}
\vspace*{-2mm}
\begin{itemize}[leftmargin=0.6cm, label=\textbullet]
	\setlength\itemsep{0.2em}
	\item \textbf{On Infinite Tensor Networks, Complementary Recovery and Type II Factors}. W. Chemissany, E. Gesteau, A. Jahn, D. Murphy, \textbf{L. Shaposhnik}. J. Phys. A 58 (2025) 435301.
	\item \textbf{Searching for Emergent Spacetime in Spin Glasses}. D. Saraidaris, \textbf{L. Shaposhnik}. JHEP 04 (2026) 032.
\end{itemize}

\section{Languages}
\vspace*{-2mm}
\setlength{\columnsep}{8pt} % tighten spacing between the 4 columns (default is larger)
\begin{multicols}{4}
\begin{itemize}[label=\textbullet, leftmargin=*, nosep, topsep=0pt, partopsep=0pt, parsep=0pt, itemsep=0pt]
	\item \textbf{English} (Fluent)
	\item \textbf{German} (Native)
	\item \textbf{Russian} (Native)
	\item \textbf{Ukrainian} (Beginner)
\end{itemize}
\end{multicols}

\end{document}
